\section{Los cambios de paradigma en los grandes modelos globales durante el último año}

\subsection{Primer paradigma: modelos de razonamiento con pensamiento intensivo}

Los modelos de razonamiento (Reasoning Model) basados en aprendizaje por refuerzo, representados por O1, son uno de los avances técnicos más importantes del último año.

\begin{importantbox}{La esencia de los modelos de razonamiento: el ciclo de conjetura y verificación}
La capacidad central de los modelos de razonamiento es la «reflexión» (Reflection), que en esencia comprende dos capacidades:
\begin{enumerate}
    \item \textbf{Formular conjeturas}: proponer continuamente nuevas hipótesis durante el proceso de resolución del problema
    \item \textbf{Autoverificación}: verificar de manera implícita cada conjetura para determinar si es correcta o incorrecta
\end{enumerate}
Mediante múltiples ciclos de «conjetura → verificación → nueva conjetura → nueva verificación», el modelo puede convertir un problema de Pass@2 en uno de Pass@1: en esencia, es equivalente a haberlo intentado varias veces. Esto se parece mucho al proceso humano de investigación científica y resolución de problemas.
\end{importantbox}

Yang Zhilin (杨植麟) compara de manera gráfica este tipo de modelo de razonamiento con un \textbf{«cerebro en una cubeta»}: no necesita interactuar con el mundo exterior, sino que solo piensa dentro de su propio cerebro, y así puede resolver problemas mediante el ciclo de conjetura y verificación.

\begin{knowledgebox}{Muestreo serial vs. paralelo}
La forma efectiva de trabajo de los modelos de razonamiento es principalmente serial: cada conjetura se basa en el resultado de la anterior. Aunque también se puede combinar con una estrategia paralela (muestrear varias soluciones al mismo tiempo), las investigaciones recientes muestran que el límite superior de la estrategia serial suele ser más alto. Esto coincide con las conclusiones experimentales de Moonshot AI (月之暗面).
\end{knowledgebox}

\subsection{Segundo paradigma: el aprendizaje por refuerzo de Agent}

A diferencia del «cerebro en una cubeta» de los modelos de razonamiento, el modelo de Agent interactúa con el entorno exterior mediante múltiples turnos:

\begin{itemize}
    \item Piensa y ejecuta acciones al mismo tiempo (búsqueda, navegador, escritura de código, etc.)
    \item Resuelve problemas mediante un método de múltiples turnos
    \item La siguiente acción depende de la retroalimentación externa y de la actualización del estado
\end{itemize}

\begin{importantbox}{La unificación de dos formas de Test-Time Scaling}
Los modelos de razonamiento y los modelos de Agent apuntan al mismo concepto central: \textbf{Test-Time Scaling}:
\begin{itemize}
    \item \textbf{Modelos de razonamiento}: hacen Scale de la cantidad de tokens de pensamiento dentro de un solo turno (más razonamiento interno)
    \item \textbf{Modelos de Agent}: hacen Scale del número de turnos de interacción (más interacción externa)
\end{itemize}
Ambos consisten en invertir más recursos de cómputo en el momento de la inferencia para completar tareas más complejas. Un modelo de conversación tradicional puede generar solo unos cientos de tokens, mientras que un modelo de razonamiento o de Agent puede tardar varias horas y consumir una gran cantidad de tokens para completar una tarea compleja.
\end{importantbox}

\subsection{Claude vs. OpenAI: diferencias de ruta técnica}

\begin{knowledgebox}{Reasoning y Agent no tienen necesariamente una dependencia lineal}
Yang Zhilin señala una observación interesante: los modelos de Claude no destacan en los benchmarks de Reasoning, pero tienen un desempeño muy bueno en las tareas de Agent. Esto indica que:
\begin{itemize}
    \item Reasoning y Agent corresponden a dos dimensiones distintas de Test-Time Scaling
    \item En el desarrollo, no existe un orden secuencial necesario entre ambos
    \item Pero para alcanzar el nivel más alto de Agent, al final sigue siendo necesaria una capacidad de Reasoning muy sólida
\end{itemize}
Es posible que Claude haya invertido más en la dirección de Agent (interacción de múltiples turnos), mientras que OpenAI ha invertido más en la dirección de Reasoning (pensamiento profundo). Las distintas rutas técnicas generan diferencias de producto en el corto plazo, pero a largo plazo ambas son indispensables.
\end{knowledgebox}

\subsection{Tercera tendencia: las empresas de modelos hacen producto propio}

\begin{importantbox}{Diseño directo vs. ingeniería inversa}
Yang Zhilin distingue dos formas de desarrollar productos:
\begin{itemize}
    \item \textbf{Ingeniería inversa de terceros}: a partir de un modelo base existente, se hace ingeniería inversa del proceso de entrenamiento del modelo mediante el diseño de andamiajes, herramientas y System Prompt, para encontrar la mejor forma de usarlo (como Manus, entre otros)
    \item \textbf{Diseño directo propio}: la empresa de modelos primero diseña las herramientas y el método de Context Engineering, y después entrena el modelo en ese entorno, de modo que el modelo rinda mejor de forma natural en su propio entorno (como Claude Code, ChatGPT Agent)
\end{itemize}
La segunda forma probablemente tiene un límite superior más alto, porque permite integrar mejor las herramientas y el modelo, y entrenar de extremo a extremo.
\end{importantbox}

\subsection{Resumen de la sección}

Los tres grandes cambios de paradigma del último año son: (1) los modelos de razonamiento (el ciclo de conjetura y verificación al estilo del cerebro en una cubeta), (2) los modelos de Agent (múltiples turnos de interacción con el mundo exterior) y (3) las empresas de modelos que hacen producto propio. El razonamiento y Agent son dos dimensiones de Test-Time Scaling que, al final, deben combinarse.

